% !TEX program = pdflatex
% ------------------------------------------------------------
% MATH 521 Homework Template (adapted from your MATH 421 HW10)
%
% HOW TO USE (quick):
%   1) Edit the Metadata block (HW number / due date especially).
%   2) Paste each problem statement under \problem[]{<n>}.
%   3) Write in the Scratch/Discussion box if you want.
%   4) Write your proof under "Proof." and end with \qed (or \qedhere).
%
% Notes:
% - Keeps theorem/definition boxes (tcolorbox).
% - Keeps ToC + PDF bookmarks for each problem.
% - Adds a Scratch/Discussion box per problem (optional).
% - Header bar has course + meeting info (left) and HW number (right).
% ------------------------------------------------------------

% TROUBLESHOOTING NOTE:
% If you ever see errors like `\tcb@insert@after@part` or other tcolorbox internals being
% an "Undefined control sequence", that usually means LaTeX is accidentally loading a
% DIFFERENT (often older) `tcolorbox.sty` from your project folder or local texmf tree.
% Fix: search your repo for `tcolorbox.sty` and delete/rename any local copy so TeX uses
% the TeX Live version. Also try "Clean" in LaTeX Workshop to remove stale .aux/.toc/.out.

\documentclass[12pt]{article}

% =========================
% 0) Metadata (EDIT ME)
% =========================
\newcommand{\CourseCode}{MATH 521}
\newcommand{\CourseTitle}{Analysis I}
\newcommand{\CourseSection}{LEC 002}
\newcommand{\Semester}{Spring 2026}
\newcommand{\Instructor}{Thierry Laurens}
\newcommand{\MeetingInfo}{MWF 1:20--2:10 \;|\; VV B119} % optional
\newcommand{\HomeworkNumber}{8} % <-- change each week
\newcommand{\YourName}{Alejandro De La Torre}
\newcommand{\DueDate}{March 20, 2026} % <-- or set e.g. January 31, 2026

% =========================
% 1) Core math + proof tools
% =========================
\usepackage{amsmath, amssymb, amsthm}

% =========================
% 2) Lists and spacing
% =========================
\usepackage{enumitem}
\setlist[enumerate,1]{label=\textbf{(\alph*)}, leftmargin=2em, itemsep=0.25em}
\setlist[itemize]{leftmargin=1.5em, itemsep=0.25em}
\setlength{\parskip}{0.35em}
\setlength{\parindent}{0pt}

% =========================
% 3) Page geometry + header/footer
% =========================
\usepackage{geometry}
\geometry{margin=1in}

\usepackage{fancyhdr}
\pagestyle{fancy}
\fancyhf{}
\lhead{\CourseCode\ --- \CourseSection, \Semester \;(\MeetingInfo)}
\rhead{Homework \HomeworkNumber}
\cfoot{\thepage}
\renewcommand{\headrulewidth}{0.4pt}
% Fixes common fancyhdr warning about header height:
\setlength{\headheight}{15pt}
% Optional: reclaim a bit of vertical space after increasing headheight
\addtolength{\topmargin}{-2pt}

% =========================
% 4) Links (subtle)
% =========================
\usepackage[hidelinks]{hyperref}
\setcounter{tocdepth}{1}

% =========================
% 5) Quotes / figures
% =========================
\usepackage{csquotes}
\usepackage{graphicx}
\graphicspath{{images/}} % put figures in ./images/

% =========================
% 6) Simple scratch box (no tcolorbox)
% =========================
% A lightweight environment for brief planning / scratch work.
% This avoids any tcolorbox dependencies.
\newenvironment{scratch}{%
  \par\smallskip
  \noindent\textbf{Discussion \& Scratch Work}\begin{quote}\small
}{%
  \end{quote}\smallskip
}

% =========================
% 7) Lightweight theorem-like environments (optional)
% =========================
\newtheorem*{claim}{Claim}
\newtheorem*{lemma}{Lemma}
\newtheorem*{remark}{Remark}
\newtheorem{theorem}{Theorem}
\newtheorem{proposition}{Proposition}
\newtheorem{corollary}{Corollary}
\newtheorem{definition}{Definition}
\newtheorem{example}{Example}
\newtheorem{exercise}{Exercise}

% =========================
% 8) Problem header command (with TOC + PDF bookmarks)
% Usage: \problem[Optional short title]{<number>}
% =========================
\makeatletter
\newcommand{\problem}[2][]{%
  \def\prob@title{Problem #2\if\relax\detokenize{#1}\relax\else\ (\textit{#1})\fi}%
  \phantomsection%
  \pdfbookmark[1]{\prob@title}{prob#2}%
  \addcontentsline{toc}{section}{\prob@title}%
  \section*{\prob@title}%
  \vspace{-0.25em}\hrule\vspace{0.8em}%
}
\makeatother

% =========================
% 9) Title block
% =========================
\title{Homework \HomeworkNumber}
\author{\YourName \\
\CourseCode\ --- \CourseTitle \\
\CourseSection \\
Instructor: \Instructor}
\date{\DueDate}

\begin{document}
\maketitle

% ------------------ collaboration + sources ------------------
% \section*{Collaboration Statement}
% Write “I worked alone” or list collaborators / external resources.
% "I worked on this homework independently. No outside collaborators or resources were used."

\tableofcontents
\bigskip
\hrule
\bigskip

\newpage

% ============================================================
% Problems (duplicate blocks as needed)
% ============================================================

\problem[]{1}
Consider the metric space $(\mathbb{R}, d)$ with the usual metric
$d(x,y)=|x-y|$. For each of the following sets $A\subseteq \mathbb{R}$, find
its interior and prove your answer.
\begin{enumerate}[label=\textbf{(\alph*)}]
  \item $A=[0,1)$
  \item $A=\mathbb{Z}$
  \item $A=\{x\in\mathbb{Q}: 0<x<1\}$
\end{enumerate}

\begin{scratch}
% Outline the strategy, key definitions, lemmas, or a proof sketch.
\end{scratch}

\textbf{Proof.}\\
% Write your proof here.
\qed

\newpage

\problem[]{2}
Recall that $d_1$ and $d_\infty$ are both metrics on $\mathbb{R}^2$.
\begin{enumerate}[label=\textbf{(\alph*)}]
  \item Prove that $d_\infty(x,y)\le d_1(x,y)\le 2d_\infty(x,y)$ for all
  $x,y\in\mathbb{R}^2$.
  \item Prove that for any $x\in\mathbb{R}^2$ and $r>0$ we have
  \[
  \{y\in\mathbb{R}^2:d_\infty(x,y)<\tfrac{r}{2}\}\subseteq
  \{y\in\mathbb{R}^2:d_1(x,y)<r\}\subseteq
  \{y\in\mathbb{R}^2:d_\infty(x,y)<r\}.
  \]
  I recommend that you draw a picture of these three sets.
  \item Prove that a set $A\subseteq \mathbb{R}^2$ is open in the metric space
  $(\mathbb{R}^2,d_1)$ if and only if it is open in the metric space
  $(\mathbb{R}^2,d_\infty)$.
\end{enumerate}

\begin{scratch}
\end{scratch}

\textbf{Proof.}\\
\qed

\newpage

\problem[]{3}
Let $X$ be a nonempty set, and recall that the discrete metric $d$ is a metric
on $X$. Prove that every set $A\subseteq X$ is open in the metric space
$(X,d)$.

\begin{scratch}
\end{scratch}

\textbf{Proof.}\\
\qed

\newpage

\problem[]{4}
Let $(X,d)$ be any metric space. Define the function $\tilde{d}:X\times X\to
\mathbb{R}$ by
\[
\tilde{d}(x,y)=\min\{d(x,y),1\}.
\]
\begin{enumerate}[label=\textbf{(\alph*)}]
  \item Prove that $\tilde{d}$ is also a metric on $X$.
  \item Prove that a set $A\subseteq X$ is open in $(X,d)$ if and only if it is
  open in $(X,\tilde{d})$.
\end{enumerate}

\begin{scratch}
\end{scratch}

\textbf{Proof.}\\
\qed

\newpage

\problem[]{5}
Let $(X,d)$ be a metric space, $A\subseteq X$, and $x\in X$. Prove that $x$ is
in $A^\circ$ if and only if there exists an open set $U$ so that
$x\in U\subseteq A$.

\begin{scratch}
\end{scratch}

\textbf{Proof.}\\
\qed

\newpage

\problem[]{6}
Let $(X,d)$ be a metric space and let $A$ be a nonempty subset of $X$. Prove
that $A$ is open if and only if there exists a set $I$, positive numbers
$r_i>0$, and points $a_i\in X$ so that
\[
A=\bigcup_{i\in I} B_{r_i}(a_i).
\]

\begin{scratch}
\end{scratch}

\textbf{Proof.}\\
\qed

% ============================================================
% Lecture 20-22 Material Proven In Class
% ============================================================
\newpage
\appendix
\section*{Appendix: Toolbox}
\addcontentsline{toc}{section}{Appendix: Toolbox (optional)}

\subsection*{Metrics and Metric Spaces}
\begin{definition}
Let $X$ be a nonempty set. A metric on $X$ is a function
$d:X\times X\to\mathbb{R}$ such that
\begin{enumerate}[label=\textbf{(\arabic*)}]
  \item $d(x,y)\ge 0$ for all $x,y\in X$.
  \item $d(x,y)=0$ if and only if $x=y$.
  \item $d(x,y)=d(y,x)$ for all $x,y\in X$.
  \item $d(x,z)\le d(x,y)+d(y,z)$ for all $x,y,z\in X$.
\end{enumerate}
\end{definition}

\begin{definition}
If $d$ is a metric on $X$, then $(X,d)$ is called a metric space.
\end{definition}

\subsection*{$\ell^p$ Metrics on $\mathbb{R}^n$}
\begin{proposition}[H\"older's inequality]
Fix $1\le p\le\infty$, and let $1\le q\le\infty$ satisfy
\[
\frac{1}{p}+\frac{1}{q}=1.
\]
For any $(x_1,\dots,x_n),(y_1,\dots,y_n)\in\mathbb{R}^n$ we have
\[
\sum_{k=1}^n |x_k y_k|
\le
\left(\sum_{k=1}^n |x_k|^p\right)^{1/p}
\left(\sum_{k=1}^n |y_k|^q\right)^{1/q}.
\]
When $p=\infty$, replace $\left(\sum_{k=1}^n |x_k|^p\right)^{1/p}$ by
$\max_{k=1,\dots,n} |x_k|$ above.
\end{proposition}

\begin{remark}
When $p=2$ (and so $q=2$), we call this the Cauchy--Schwarz inequality.
\end{remark}

\begin{corollary}[Minkowski's inequality]
For any $1\le p\le\infty$ and $(x_1,\dots,x_n),(y_1,\dots,y_n)\in\mathbb{R}^n$,
we have
\[
\left(\sum_{k=1}^n |x_k+y_k|^p\right)^{1/p}
\le
\left(\sum_{k=1}^n |x_k|^p\right)^{1/p}
+
\left(\sum_{k=1}^n |y_k|^p\right)^{1/p}.
\]
\end{corollary}

\begin{corollary}
For any $1\le p\le\infty$, the function $d_p:\mathbb{R}^n\times\mathbb{R}^n\to\mathbb{R}$,
\[
d_p(x,y)=
\begin{cases}
\left(\sum_{k=1}^n |x_k-y_k|^p\right)^{1/p} & \text{if } 1\le p<\infty,\\
\max_{k=1,\dots,n} |x_k-y_k| & \text{if } p=\infty,
\end{cases}
\]
is a metric on $\mathbb{R}^n$.
\end{corollary}

\subsection*{Open Balls and Open Sets}
\begin{definition}
Let $(X,d)$ be a metric space. The open ball (or neighborhood) of radius $r>0$
centered at $a\in X$ is
\[
B_r(a)=\{x\in X:d(a,x)<r\}.
\]
\end{definition}

\begin{definition}
Let $(X,d)$ be a metric space and $A\subseteq X$. We say:
\begin{enumerate}[label=\textbf{(\arabic*)}]
  \item $a\in X$ is an interior point of $A$ if there exists $r>0$ such that
  $B_r(a)\subseteq A$.
  \item
  \[
  A^\circ=\{a\in X: a \text{ is an interior point of } A\}
  \]
  is the interior of $A$.
  \item $A$ is open if $A^\circ=A$.
\end{enumerate}
\end{definition}

\begin{remark}
Note that $A^\circ\subseteq A$ for any set $A$.
\end{remark}

\subsection*{Basic Properties of Interior and Open Sets}
\begin{proposition}
Let $(X,d)$ be a metric space and $A,B\subseteq X$.
\begin{enumerate}[label=\textbf{(\arabic*)}]
  \item If $A\subseteq B$, then $A^\circ\subseteq B^\circ$.
  \item $A^\circ\cup B^\circ\subseteq (A\cup B)^\circ$.
  \item $A^\circ\cap B^\circ=(A\cap B)^\circ$.
  \item $(A^\circ)^\circ=A^\circ$. In particular, $A^\circ$ is an open set.
  \item $A^\circ$ is the largest open set contained in $A$.
  \item A finite intersection of open sets is open.
  \item An arbitrary union of open sets is open.
\end{enumerate}
\end{proposition}

\begin{remark}
An arbitrary intersection of open sets need not be open.
\end{remark}

\end{document}
